%======================================================================
\section{Prologue: One Theorem, Ten Perspectives}
%======================================================================

\subsection*{New Year's Night, 1802}

On the first day of the nineteenth century, the astronomer Giuseppe Piazzi
spotted something that should not have been there: a faint, slow-moving
point of light in the constellation Taurus. Night after night he tracked
it, recording its position against the fixed stars.  He had discovered
Ceres---the first asteroid, and, as it would turn out, the largest object
between Mars and Jupiter.

Then disaster struck.  After only 41 nights of observation, Ceres
drifted too close to the Sun in the sky and vanished in the glare.
Piazzi had captured barely one percent of the asteroid's orbit.  As the
weeks became months, the astronomical community faced a crisis: by the
time Ceres emerged from behind the Sun, it could be \emph{anywhere} along
a vast arc of sky.  The needle had disappeared into a haystack of stars,
and no one knew where to point a telescope.

No one, that is, except a 24-year-old mathematician in Brunswick named
Carl Friedrich Gauss.

Gauss took Piazzi's 41 data points---each one contaminated by
atmospheric distortion, instrument imprecision, and the tremor of a
human hand---and asked a question that would reshape the mathematics of
uncertainty: \emph{Of all the elliptical orbits that the laws of physics
permit, which one passes closest to these noisy observations?}

His method was, at its core, breathtakingly simple.  He wrote down the
equations of a Keplerian ellipse, parameterised by six orbital elements.
These equations carved out a surface---a subspace of possible
trajectories---inside the much larger space of all conceivable sequences
of celestial positions.  Piazzi's 41 observations lived somewhere in that
larger space, displaced from the true orbit by the accumulated noise of
every measurement.  Gauss sought the single point on the orbital surface
that lay \emph{nearest} to the data.

He found it by dropping a perpendicular.

The computation was ferocious: weeks of hand calculation, pages of
logarithm tables, corrections for the Sun's gravitational tug.  But the
geometric principle was crystalline.  The best-fitting orbit was the
\textbf{shadow} of the noisy data cast onto the subspace of physically
possible trajectories---the point where a perpendicular from the data
meets the surface.  The residual---the gap between data and
prediction---stood at a right angle to every direction the orbit could
vary.

Gauss sent his prediction to the astronomer Franz Xaver von Zach. On
December~31, 1801---New Year's Eve---von~Zach aimed his telescope at the
patch of sky Gauss had specified. The clouds parted.  And there, within
half a degree of the predicted position, was Ceres.

It was one of the greatest predictions in the history of science. And it
was, at its mathematical heart, an orthogonal projection.

\bigskip

\subsection*{From the stars to the spreadsheet}

What Gauss did for planetary orbits, you do every time you run a
regression.  The setting is humbler---you might be predicting house
prices instead of asteroid trajectories---but the geometry is identical.
Your data vector $\by\in\R^n$ lives in a high-dimensional space.  The
column space of your design matrix $\Col(\bX)$ is a subspace of
possible fitted values, one for each choice of coefficients.  The
ordinary least squares fit $\yhat$ is the point in $\Col(\bX)$ closest
to $\by$.  The residual $\be = \by - \yhat$ is perpendicular to the
column space.

Gauss projected noisy observations onto the subspace of Keplerian
ellipses.  You project noisy responses onto the subspace spanned by
your regressors.  The theorem that makes both possible is the same.

\begin{metaphor}
Imagine a flashlight shining straight down onto a flat wall.
The \textbf{object} is the data vector $\by\in\R^n$.
The \textbf{wall} is the column space $\Col(\bX)$.
The \textbf{shadow} is the fitted vector $\yhat$.
The \textbf{gap} between object and shadow is the residual $\be$.
The shadow always falls at a \textbf{right angle}---the residual is
perpendicular to the wall.
\end{metaphor}

This single image---a shadow cast at right angles---generates everything
in these notes.

\begin{theorem}[The Projection Theorem]\label{thm:projection}
Let $S$ be a closed subspace of an inner product space $V$, and let
$\by\in V$. There exists a unique $\yhat\in S$ such that
\[
  \norm{\by - \yhat} \;\le\; \norm{\by - \bm{s}}
  \qquad\text{for all } \bm{s}\in S.
\]
Moreover, $\yhat$ is characterised by the \emph{orthogonality condition}
\[
  (\by - \yhat) \perp S,
\]
i.e.\ $\ip{\by - \yhat}{\bm{s}} = 0$ for every $\bm{s}\in S$.
\end{theorem}

Ordinary least squares is the Projection Theorem applied to
$V = \R^n$ with the Euclidean inner product and $S = \Col(\bX)$.
That one sentence is the seed from which the rest of these notes grow.

\subsection*{Roadmap}

Every section that follows answers a question that the Projection
Theorem makes natural:

\begin{enumerate}[leftmargin=2em]
  \item \textbf{Where do the coefficients come from?}
    The OLS coefficients $\bbetahat$ are the \emph{coordinates} of the
    shadow in the basis provided by the columns of $\bX$.
    (\S\ref{sec:coefficients})

  \item \textbf{How good is the fit?}
    The Pythagorean theorem decomposes total variance into explained
    and unexplained parts, and $R^2$ turns out to equal the squared
    cosine of the angle between $\by$ and $\yhat$.
    (\S\ref{sec:pythagoras})

  \item \textbf{What does each variable contribute on its own?}
    The Frisch--Waugh--Lovell theorem shows that the coefficient on any
    single regressor comes from projecting away all the others
    first---a geometric ``peeling apart'' of contributions.
    (\S\ref{sec:fwl})

  \item \textbf{Which data points matter most?}
    Leverage measures how far each observation sits from the centre of
    the regressor cloud; influence measures how much the shadow shifts
    when you remove a single point.
    (\S\ref{sec:leverage})

  \item \textbf{What can go wrong?}
    Multicollinearity collapses the subspace toward a thin sliver,
    making the shadow unstable.  Omitted variables tilt the subspace
    itself, so the shadow falls in the wrong place.
    (\S\ref{sec:failures})

  \item \textbf{Can we do better by accepting some bias?}
    Ridge and LASSO regularisation deliberately \emph{shrink} the
    shadow away from the perpendicular foot, trading a little bias for
    a lot less variance.
    (\S\ref{sec:regularisation})

  \item \textbf{How do we test whether the shadow is real?}
    The $F$-test compares the length of two shadows---one cast onto the
    full column space, the other onto a restricted subspace---and asks
    whether the difference is too large to be noise.
    (\S\ref{sec:ftest})
\end{enumerate}

\noindent
By the time you reach the end, you will see that Gauss's prediction of
Ceres, the $R^2$ on your next problem set, and the $p$-value in a
clinical trial are all manifestations of the same geometric act:
dropping a perpendicular onto a subspace, and measuring the gap that
remains.

Let us begin.
